% SGA 3, Expose V, Loop 1 English reconstruction.
% Scope: Lemma 7.2, including its closing seam on component-local p. 26.
% Authority: Polo--Gille Exp5-13oct24.pdf, component-local pp. 25--26
% (combined-reader pp. 287--288), re-edited headers 275--276.
% Jacob Reinhold's CC BY 4.0 Markdown is comparison/drafting material only.

\medskip
\phantomsection\label{lem:sga3-V-7.2}%
\noindent\textbf{引理 7.2.}---%
\emph{设 \(T\) 是仿射 Noether 概形, \(X,Y,Z\) 是有限型
\(T\)-概形, 且 \(X\) 在 \(T\) 上拟射影, 并设}
\SGAThreeDiagram[0.30\textwidth]
  {figures/exp5/p275_lemma72_square.png}
  {\hyperref[lem:sga3-V-7.2]{引理~7.2} 中的图.}
  {fig:sga3-V-7.2-square}
  {第 V 篇, 分篇第~25 页 / 合订本第~287 页 / 重版本页眉页码~275}
\emph{是 \((\mathbf{Sch}/T)\) 中的图.
设 \(z\) 是 \(v(Y)\) 中相对于 \(T\) 闭的点,
并假设 \(v\) 在 \(v^{-1}(z)\) 的各点处平坦.
则存在 \(X\) 的一个闭子概形 \(F\), 使得
\[
  u\bigl(u^{-1}(F)\cap v^{-1}(z)\bigr)
\]
有限且非空, 并且 \(v\) 在 \(u^{-1}(F)\) 上的限制
在 \(v^{-1}(z)\) 的各点处平坦.}

置 \(T=\operatorname{Spec}A\). 可假设
\[
  X=\operatorname{Proj}\mathscr S,
\]
其中 \(\mathscr S\) 是有限生成 \(A\)-模 \(E\) 的对称代数.

若 \(u(v^{-1}(z))\) 有限, 可取 \(F=X\). 否则, 设
\[
  y_1,\ldots,y_n
\]
是纤维 \(v^{-1}(z)\) 的结构层 \(\mathcal O_{v^{-1}(z)}\) 的相伴点.
于是, 若 \(\mathcal O_i\) 表示 \(v^{-1}(z)\)
在 \(y_i\) 处的局部环, 则 \(\mathcal O_i\) 的极大理想由零因子组成.
设 \(t\) 是 \(z\) 在 \(T\) 中的像.
集合 \(u(v^{-1}(z))\) 是 \(X\) 在 \(t\) 上的纤维中的无限可构造子集.
因此, 该纤维包含一个闭点 \(x\in u(v^{-1}(z))\),
它异于 \(u(y_1),\ldots,u(y_n)\).
开集 \(X-\{x\}\) 是 \(u(y_1),\ldots,u(y_n)\) 的邻域,
因而包含一个形如 \(D_+(f)\) 的开邻域,
其中 \(f\in\mathscr S\) 为 \(d\) 次齐次元
(使用 EGA~II, \S2.3 的记号).

因此, 由 \(f\) 定义的闭子概形
\[
  X_1=V_+(f)
\]
包含 \(x\) 且避开各点 \(u(y_1),\ldots,u(y_n)\).
其逆像
\[
  Y_1=u^{-1}\bigl(V_+(f)\bigr)
\]
因而异于 \(Y\), 且与 \(v^{-1}(z)\) 相交.
事实上, 我们将证明 \(v\) 在 \(Y_1\) 上的限制 \(v_1\)
在 \(v^{-1}(z)\) 的各点处平坦.
若 \(u(v_1^{-1}(z))\) 有限, 就可取 \(F=X_1\).
否则, 用 \(Y_1,v_1\) 以及 \(u\) 在 \(Y_1\) 上诱导的态射
\(u_1\) 替代 \(Y,v,u\), 重复这一论证.
这样就得到 \(X\) 的闭子概形降链
\[
  X\supset X_1\supset\cdots
\]
由于这样的降链会终止, 对某个 \(n\),
\[
  u\bigl(u^{-1}(X_n)\cap v^{-1}(z)\bigr)
\]
有限且非空, 此时取 \(F=X_n\).

还须证明 \(v_1\) 在 \(v^{-1}(z)\) 的各点处平坦.
设 \(y\in Y_1\) 位于 \(z\) 上,
\(\mathcal O_y\) 是 \(Y\) 在 \(y\) 处的局部环,
\(\overline{\mathcal O}_y\) 是 \(v^{-1}(z)\) 在 \(y\) 处的局部环,
\(\mathcal O_{v(y)}\) 是 \(Z\) 在 \(v(y)\) 处的局部环.
选取 \(g\in\mathscr S_1\), 使 \(D_+(g)\) 是 \(u(y)\) 在 \(X\) 中的邻域,
并设 \(\varphi\) 和 \(\overline\varphi\)
分别是 \(f/g^d\) 在 \(\mathcal O_y\) 和 \(\overline{\mathcal O}_y\) 中的像.
由 \(f\) 的构造, 元素 \(\overline\varphi\)
不是 \(\overline{\mathcal O}_y\) 中的零因子.
由于 \(\mathcal O_y\) 在 \(\mathcal O_z=\mathcal O_{v(y)}\) 上平坦,
元素 \(\varphi\) 不是 \(\mathcal O_y\) 中的零因子, 且
\[
  \mathcal O_y/\mathcal O_y\varphi
\]
在 \(\mathcal O_z\) 上平坦 (SGA~1, IV, 5.7).
此商恰是 \(Y_1\) 在 \(y\) 处的局部环.
% SGA3-V-V2-EOL-PRESERVATION

% Translation coverage: entire assigned component; adjoining components retain the original statement/proof seams.
% Source lines: 96; translated body lines before this audit comment: 91.
% Math preserved: 92 inline formulas, 8 display blocks, 0 equation environments; mathematical prose translated where present.
% Labels: 1/1; hyperrefs: 1/1; SGARefLinks: 0/0; inline targets: 0/0; footnotes: 0/0.
% Figure arguments preserved; all rendered captions and natural-language source locators translated.
% Unresolved TODOs: 0; compilation and visual inspection remain with the volume integrator.
